\documentclass[11pt]{article}
\usepackage[OT1]{fontenc}
\usepackage[margin=1in]{geometry}
\usepackage{amsmath,amssymb,booktabs,longtable,array,calc,hyperref}
\hypersetup{colorlinks=true,linkcolor=blue,urlcolor=blue}
\providecommand{\pandocbounded}[1]{#1}
\providecommand{\tightlist}{}
\title{A finiteness theorem for elementary-positive three-legged spiders}
\author{Leo Gao\\\small AI-assisted research}
\date{26 November 2026}
\begin{document}\maketitle


This
is a partial theorem about Tom's classification conjecture, not a claim
to have proved its list of seven graphs.

\textbf{Theorem.} For every fixed integer \(c\ge2\) other than 3, only
finitely many three-legged spiders \(S(a,b,c)\), with \(a\ge b\ge c\),
have elementary-positive chromatic symmetric function. In particular
this proves the finiteness assertion for the entire previously unbounded
family \(S(a,b,2)\). For \(c=3\), any infinite counterfamily to
finiteness must, after dropping finitely many terms, escape along the
balanced direction \(a/b\to1\), and Section 5 sharpens this to two
explicit arithmetic progressions: \(a=b+14\) and
\(b\equiv16,36\pmod{140}\).

\emph{Terminology.} A spider has a common vertex joined to paths of
edge-lengths \(a,b,c\); it has \(n=a+b+c+1\) vertices.
Elementary-positive means that every coefficient of its chromatic
symmetric function in the elementary basis is nonnegative.

\section*{1. A randomized coefficient identity
(exact)}\label{a-randomized-coefficient-identity-exact}

Write \(P_j=X_{\mathrm{path}_j}\) and \(d=c+1\). We use two established
identities: the path expansion \[
P_n=\sum_{I\vDash n}i_1\prod_{t>1}(i_t-1)e_I
\] and the spider identity \[
X_S=P_n+\sum_{i=1}^c(P_iP_{n-i}-P_{b+i}P_{n-b-i}).
\] See \href{https://arxiv.org/pdf/2112.06619}{Wang--Wang, Proposition
2.5 and Lemma 2.10} (the latter attributed to Zheng).

Choose a partition consisting of \(M-L\) copies of \(d\) and \(L\)
copies of \(d+1\), so \(dM+L=n\). Let \(V_1,\ldots,V_M\) be these parts
in uniformly random order. Write \(T_j=\sum_{i\le j}V_i\),
\(R(x)=x/(x-1)\), and \(H=\{T_j\in[b+1,b+d-1]\text{ for some }j\}\). On
\(H\) this \(j\) is unique; write it \(\tau\). The selected coefficient
divided by \(\binom ML(d-1)^{M-L}d^L>0\) is exactly \[
 C=\mathbb E R(V_1)-\mathbb E[\mathbf1_H R(V_1)R(V_{\tau+1})].
\] Indeed no part fits in the small positive factors \(P_i\), and each
negative product is obtained by cutting an ordered composition at the
indicated partial sum, introducing exactly one additional first-part
weight. Both segments are nonempty. Normalize by the number of orderings
and the product of all part-minus-one factors.

\section*{2. A sparse-defect limit}\label{a-sparse-defect-limit}

Suppose \(b\to\infty\). Pass to a subsequence on which
\(p=\lim b/n\in[0,1/2]\) exists, and the residues \(r=b\bmod d\) and
\(s=n\bmod d\) are fixed. For \(p>0\), take any fixed nonnegative
\(L\equiv s\pmod d\). For \(p=0\), fix \(\mu>0\) and choose
\(L\equiv s\pmod d\) nearest to \(\mu n/b\). In both cases \(M=(n-L)/d\)
is integral, \(0\le L\le M\) eventually and \(L/M\to0\).

Put \(m=\lfloor b/d\rfloor\), and let \(J\) be the number of \(d+1\)'s
in the first \(m\) positions. It has the hypergeometric distribution. If
\(p>0\) then \(J\Rightarrow Z\sim\mathrm{Bin}(L,p)\), by sampling the
fixed number \(L\) of marked positions without replacement, as
\(m/M\to p\). If \(p=0\) then \(m,L\to\infty\), \(mL/M\to\mu\), and
\(J\Rightarrow Z\sim\mathrm{Pois}(\mu)\): for every fixed \(k\),
\(\mathbb E(J)_k=(m)_k(L)_k/(M)_k\to\mu^k\) in falling-factorial
notation; the Poisson law is moment-determinate.

We claim \(\Pr(H)\to1-\Pr(Z\equiv r\pmod d)\). For any fixed \(K\), the
probability of a marked position in \([m-K,m+K]\) is at most
\((2K+1)L/M\to0\). On its complement, near \(m\) we have \(T_t=dt+J\).
If \(J\le K\) and \(K\ge4\), this arithmetic progression locally spans
the entire target interval and its neighbors, since \(0\le b-dm<d\).
Hence it hits the target precisely when \(J\not\equiv r\pmod d\); any
successor part is \(d\) (increase \(K\) by two if necessary). Since
\(J\) is tight, let \(K\to\infty\). Also \(\Pr(V_1\ne d)=L/M\to0\).
Boundedness of the weights in the exact identity now yields \[
 C\ \longrightarrow\ R(d)-R(d)^2[1-\Pr(Z\equiv r\pmod d)].
\] This is strictly negative exactly when \(\Pr(Z\equiv r\pmod d)<1/d\).

\section*{3. Tuning the residue bias}\label{tuning-the-residue-bias}

Let \(\omega=e^{2\pi i/d}\). For \(0<p<1/2\), put \(z_j=1-p+p\omega^j\),
\(\theta=\arg z_1\). The root-of-unity filter gives \[
 d\Pr(\mathrm{Bin}(L,p)\equiv r)-1
 =\sum_{j=1}^{d-1}\omega^{-rj}z_j^L.
\] The two conjugate terms with \(j=1,d-1\) strictly dominate in modulus
the remaining terms, since \(|z_j|^2=1-2p(1-p)(1-\cos(2\pi j/d))\).
Further \(0<\theta<\pi/d\): the chord point is strictly between the
positive real axis and the ray through the midpoint of \(1,\omega\).
Along \(L=s+kd\), the phase \(L\theta-2\pi r/d\) rotates in steps
strictly between 0 and \(\pi\). There are infinitely many \(k\) on which
its cosine is bounded above by a fixed strictly negative number: for
irrational rotation use density; for rational rotation its orbit is a
regular polygon with at least three vertices, not contained in a closed
semicircle. Choose any sufficiently large such \(L\). The sum is
strictly negative.

If \(p=1/2\), the first terms have phase \((L-2r)\pi/d\), because
\((1+\omega)/2=e^{\pi i/d}\cos(\pi/d)\). If \(\cos((s-2r)\pi/d)\ne0\),
increasing \(L\) by \(d\) flips its sign, giving the same result. If it
is zero, then \(d\) is even and \(s-2r\equiv d/2\pmod d\). For
\(d\ge6\), the first pair cancels identically, but the next pair
\(j=2,d-2\) has positive modulus \(\cos(2\pi/d)^L\) and cosine
\(\cos(2(L-2r)\pi/d)=-1\). Its magnitude strictly dominates all
remaining terms for large \(L\), giving a negative sum. Only \(d=4\) is
exceptional (the second harmonic is zero).

If \(p=0\), the filter instead gives
\(\sum_{j=1}^{d-1}\omega^{-rj}\exp(\mu(\omega^j-1))\). Again the first
pair dominates as \(\mu\to\infty\). Choose arbitrarily large \(\mu>0\)
satisfying \(\mu\sin(2\pi/d)-2\pi r/d\equiv\pi\pmod{2\pi}\); this gives
a strictly negative sum, as \(d\ge3\).

\section*{4. Finiteness}\label{finiteness}

A hypothetical infinite family with unbounded \(b\) yields a negative
elementary coefficient by Sections 1--3, outside the stated exception.
For bounded \(b\) there are only finitely many \(a\): every sufficiently
large \(n\) is the sum of copies of \(b+1,b+2\), by coprimality. For the
associated partition, the center block of a connected partition would
have to contain the whole middle and short legs, since neither can
contain a detached block of size at least \(b+1\). The center block
would then have at least \(b+c+1>b+2\) vertices. This missing connected
partition obstructs elementary positivity by Wolfgang's criterion (also
\href{https://arxiv.org/pdf/2112.06619}{Wang--Wang, Theorem 2.4}). For
\(c=3\), the exceptional congruence \(s-2r\equiv2\pmod4\) is precisely
\(a-b\equiv2\pmod4\). This proves the theorem.

\textbf{Status and limitations.}
\href{https://www.mat.univie.ac.at/~slc/wpapers/FPSAC2025/118.pdf}{Tom's
Conjecture 2.16} asks for an exact list of seven graphs for all
\(c\ge2\). \href{https://arxiv.org/pdf/2608.22184}{Wang--Wang,
arXiv:2608.22184v2, pp.~3 and 23} explicitly leaves \(3\mid b, b\ge12\)
open in the \(c=2\) slice; low multiples of 3 and the other residues are
already classified. Our compactness argument provides no practical
cutoff and does not identify the finite exceptional set. Nor does
fixed-\(c\) finiteness imply finiteness when \(c\) itself varies. No
independent peer review or priority certification is implied.

\section*{5. Sharpening the singular shortest-leg-three
case}\label{sharpening-the-singular-shortest-leg-three-case}

In fact the only possibly infinite family left for \(c=3\) is \[
   a=b+14,\qquad b\equiv16\text{ or }36\pmod{140}.
\] This is a necessary condition for escape, not a claim that its
members are positive.

To see this, restrict an escaping sequence to the exceptional regime.
Write \(a-b=4t+2\); we have \(t=o(b)\). If \(b\) is odd, the three legs
are odd and there is no perfect matching (each arm would require its own
center edge). As \(n\) is then even, this excludes the connected
partition into pairs. Suppose then \(b=4x+r\), \(r=0\) or \(2\). For the
partition with \(l\) fives and \(k=(n-5l)/4\) fours define \(W(h,0)=16\)
and for \(j\ge1\) put \[
 W(h,j)=(h+1)\cdots(h+j-1)(16h+15j).
\] The coefficient has the sign of the polynomial \[
 Q=12W(k,l)-\sum_{i=1}^3\sum_{j=0}^l\binom lj W(h,j)W(k-h,l-j),
 \quad h=(b+i-5j)/4,
\] where nonintegral or infeasible \(h\) are omitted. This is the
ordered-path weight after labeling each five and suppressing fixed
positive factors. For large \(x\) all integral \(h\) are feasible. The
exact polynomial expansions are reproduced by the accompanying script.
For \(r=2,l=2\) the polynomial is \(-16(16t+5)x-2(4t+5)(8t+9)<0\). For
\(r=0,l=6\) the \(x^6\) coefficient vanishes and the \(x^5\) coefficient
is \(-768(4t-13)\). The total degree is at most six, so the lower terms
are \(O((1+t)^2x^4)\) when \(t=o(x)\). Hence \(t\ge4\) cannot occur
infinitely often. This leaves \(r=0\), \(a-b=2,6,10,14\). Gap 2 is
excluded by the missing connected partition into two equal blocks of
size \(b+3\).

For the other gaps we use a finite, particularly simple variant of the
sparse-defect identity. Take bulk parts \(d\ge4\), and take finitely
many exceptional parts of sizes \(d+v_i\) with prescribed offsets,
allowing signed \(v_i\) as long as the exceptional part is at least 4 (a
congruent positive offset always suffices). Require
\(\sum v_i\equiv n\pmod d\). As \(b/n\to1/2\), the finitely many marked
positions fall before \(\lfloor b/d\rfloor\) asymptotically as
independent fair coin tosses. The exceptional positions avoid that
position's bounded neighborhood with probability tending to one. The
limiting chance of a cut in \([b+1,b+3]\) is therefore the fraction of
subsets of offsets whose sum modulo \(d\) lies in
\(\{b+1,b+2,b+3\}\pmod d\). If this exceeds \((d-1)/d\), the coefficient
is eventually negative. (This repeats Section 2 with an interval of
width 3 instead of \(d-1\).)

Here is the entire small certificate, with offsets written modulo \(d\);
a superscript is multiplicity, and the empty tuple means zero offsets.
Each entry is checked by at most 16 subset sums; omitted residues are
not asserted to work.

\begin{longtable}[]{@{}rrl@{}}
\toprule\noalign{}
gap & bulk \(d\) & residue \(b\bmod d\): offsets \\
\midrule\noalign{}
\endhead
\bottomrule\noalign{}
\endlastfoot
6 & 6 & 0: \(1^4\); 2: \(5^4\); 4: empty \\
10 & 5 & 0: \(1^4\); 1: \(4^4\); 2: \(3\); 3: empty; 4: \(2\) \\
14 & 5 & 0: \(1^3\); 2: \(4^3\); 3: \(4\); 4: \(1\) \\
14 & 7 & 0: \(1^4\); 3: \(6^4\); 4: \(5\); 5: empty; 6: \(2\) \\
\end{longtable}

The first row covers all even \(b\). For gap 14 only \(b\equiv1\pmod5\)
and \(b\equiv1,2\pmod7\) remain, in addition to \(4\mid b\). The Chinese
remainder theorem gives the asserted two classes.

\textbf{Reproducibility.} Run \texttt{python\ verify.py} (Python 3,
SymPy). It cross-checks the cut identity by direct finite enumeration,
verifies the two symbolic polynomials, the subset certificates and the
final residue classes. It also certifies, independently of the
asymptotic argument, nonpositivity in every previously unclassified
layer \(c=2, b=12,15,18\) for all \(a\ge b\). Beyond the two-coin
threshold the missing-partition argument applies; below it the script
finds an explicit negative coefficient for each outstanding case.

\end{document}
